% Regression test: a proof nested inside another proof, with \cite capture
% active. \pgph@installcite used to reinstall over its own capture macro in the
% inner proof, making \pgph@saved@cite the capture itself, so the first \cite
% inside the nested proof recursed forever (pdflatex hung; the harness runs
% pdflatex under timeout so a regression fails fast instead). The repeated outer
% \cite must reuse its slot without emitting a duplicate edge. The keys carry
% underscores on purpose: they must survive capture and .dot emission verbatim.
%
% Two shapes of nesting, which must not be attributed alike:
%   * a nested proof that proves no result of its own (here titled "Proof of
%     claim" but stating no claim environment): a sub-proof of the proof around
%     it, so its citation belongs to the enclosing result;
%   * a nested proof that follows a claim stated inside the proof: the claim is
%     a result in its own right, numbered and drawn, and its proof attaches to
%     it and not to the theorem being proved. The claim, being a step of the
%     proof it is stated in, gets an edge to the theorem, so that it hangs under
%     the result it was raised to establish rather than floating free.
% After the claim's proof closes, the enclosing proof must go back to citing for
% its own result: the key cited both before and after the claim keeps its single
% slot and yields one edge, to the theorem.
% Uses only amsthm; single-pass and bibtex-free (citelabel=key, manual
% \bibitem) for determinism.
\documentclass{article}
\usepackage{amsthm}
\usepackage[cite,citelabel=key]{proofgraph}
\newtheorem{theorem}{Theorem}
\newtheorem{claim}{Claim}
\begin{document}

\begin{theorem}\label{thm:main_result}Main.\end{theorem}
\begin{proof}
  Outer argument, using \cite{Outer_Ref_2011}.
  \begin{proof}[Proof of claim]
    Inner argument, using \cite[Lemma 2]{Inner_Ref_2020}.
  \end{proof}
  Outer again, with \cite{Outer_Ref_2011} reusing its slot.
\end{proof}

\begin{theorem}\label{thm:with_claim}With a claim inside its proof.\end{theorem}
\begin{proof}
  Before the claim, using \cite{Around_Ref_2013}.
  \begin{claim}\label{clm:step}A step of the argument.\end{claim}
  \begin{proof}
    The claim's own proof, using \cite{Claim_Ref_2015}.
  \end{proof}
  After the claim, using \cite{Around_Ref_2013} again.
\end{proof}

\begin{thebibliography}{9}
\bibitem{Outer_Ref_2011}An outer reference.
\bibitem{Inner_Ref_2020}An inner reference.
\bibitem{Around_Ref_2013}A reference cited around the claim.
\bibitem{Claim_Ref_2015}A reference of the claim's own proof.
\end{thebibliography}
\end{document}
